\section{Objetivo general}

Evaluar si un eDSL en TypeScript reduce la barrera de entrada a Nix ---referente funcional y declarativo en gestión de paquetes, limitado por su curva de aprendizaje--- frente a Nixlang, mediante tiempo de completación de tareas y experiencia de usuario.

\section{Objetivos específicos}

\begin{enumerate}
\item Identificar los principales puntos de dolor cognitivos que presenta el lenguaje de Nix, para fundamentar el diseño de un eDSL, mediante sesiones de pensar-en-alto y ``Programación Natural'' con estudiantes de Ciencias de la Computación que no hayan utilizado Nix previamente.
\item Desarrollar un eDSL en TypeScript que sirva de prototipo funcional para la evaluación comparativa, capaz de generar archivos de configuración en Nixlang, cubriendo al menos las funcionalidades de la librería estándar, verificado con una batería de pruebas.
\item Comparar Nixlang frente al eDSL desarrollado, para determinar si la familiaridad con TypeScript reduce la carga cognitiva de adopción mediante un cuestionario estructurado, y el uso de Short AttrakDiff 2 y DEXI aplicados a estudiantes de Ciencias de la Computación sin experiencia previa, con análisis de diferencia estadística.
\end{enumerate}
